\section{Related Work}
\label{sec:related}

\paragraph{Agents and language agents.} In artificial intelligence, an agent perceives its environment and acts on it \citep{russell2020aima}, in a cycle supported by memory. The loop that harnesses implement goes back to ReAct \citep{yao2023react}, which interleaves a model's reasoning with actions whose results return to it as observations. Later patterns live in particular primitives: Reflexion \citep{shinn2023reflexion} writes reflections on failed attempts and carries them into the next, a decision in \prim{Control flow} and \prim{Context}; MemGPT \citep{packer2023memgpt} pages information in and out of a limited window, a decision in \prim{Context}; Toolformer \citep{schick2023toolformer} trains the model to emit tool calls, moving policy into the model while the calls are still carried out by \prim{Output}.

\paragraph{CoALA.} \citet{sumers2024coala} organize language agents by memory (working, episodic, semantic, procedural), an action space of external and internal actions, and a decision cycle. CoALA describes what an agent \emph{is}, as a cognitive architecture; this paper describes what software an agent is \emph{made of}. The correspondence is in Table~\ref{tab:coala}. Two things in this paper have no counterpart in CoALA: \prim{Model interface} as a component, and the production concerns as a level of description distinct from the architecture. CoALA's memory modules, with the retrieval and learning actions that read and write them, are design decisions in \prim{Context} and \prim{Output} here, not a component.

\begin{table}[t]
\centering
\small
\caption{The five primitives, the step of the agent cycle each serves, and the CoALA component it corresponds to.}
\label{tab:coala}
\begin{tabularx}{\linewidth}{@{}lYYY@{}}
\toprule
Primitive & Mechanism & In the agent cycle & CoALA \\
\midrule
\prim{Input} & capture requests and results & perceive & observations \\
\prim{Context} & assemble what the model is shown & working and long-term memory & working memory; retrieval \\
\prim{Model interface} & request a response & decide & reasoning (by the model) \\
\prim{Output} & show words, carry out actions & act & grounding actions; learning (writes) \\
\prim{Control flow} & go again, hand back, or stop & the cycle & decision cycle \\
\bottomrule
\end{tabularx}
\end{table}

\paragraph{Agent-computer interfaces.} \citet{yang2024sweagent} show that the interface through which a model acts on a computer (the commands it is given, their output format, guardrails such as linting edits) changes performance on software tasks, and name it the agent-computer interface. In our terms the interface is a set of design decisions in \prim{Output} and \prim{Context}. The present paper asks what the complete set of components is, of which the interface is part.

\paragraph{Workflows and agents.} \citet{anthropic2024effective} distinguish workflows, in which code fixes the sequence of model calls, from agents, in which the model directs its own process, and describe patterns such as prompt chaining, routing and orchestrator--workers built on an ``augmented LLM''. We treat the distinction as a property of \prim{Control flow}, and decompose the augmented LLM into \prim{Context} (retrieval, memory) and \prim{Output} (tools) around a \prim{Model interface}.

\paragraph{The harness as a variable.} \citet{zhang2026stop} argue that harness configuration can account for more variance in performance than the choice of model, and ask that the harness be disclosed when agents are compared. \citet{fan2026empirical} hold a coding harness's loop fixed and vary planning, the action space and context management across four models, finding that each changes agent behaviour. Their three dimensions were chosen independently of this paper and each falls within one primitive: planning in \prim{Control flow}, the action space in \prim{Output}, context management in \prim{Context}. The present paper supplies what both assume: a statement of what the harness's parts are, which a disclosure can follow and a study can vary. \citet{jin2026agentprimitives} use ``primitives'' for recurring patterns of model calls in multi-agent systems (review, voting, planning and execution); in our terms these are shapes of \prim{Control flow}, and the two uses of the word are complementary.
